\scanpage{191}
% l. 01
( ceci \unsure{marche} \unsure{s.s} hyp. de finitude sur
% l. 02
$\Gamma$, le cas universel étant celui où $\Gamma$
% l. 03
est le \unsure{grpe} de Teichmüller de $\pi'$ \unsure{fixt} $i$
% l. 04
i.e. $\operatorname{Autext}_{\mathrm{lac}}(\pi';i)\simeq\operatorname{Aut}_{\mathrm{lac}}(\pi)$ ). Je
% l. 05
dis que si $\Gamma$ opère \uline{fidèlt} \uline{sur $\pi$}
% l. 06
i.e. si son opération extérieure \uline{sur}
% l. 07
\uline{$\pi'$} est fidèle, et si $\Gamma$ est \uline{fini},
% l. 08
Alors $\Gamma$ est nécessairement cyclique \inserted{(cas $\Gamma=\Gamma^0$)}
% l. 09
ou dihédral ($\Gamma\ne\Gamma^0$), et que si de
% l. 10
plus $\Gamma=\Gamma^0$, alors \deleted{\unsure{$\mathcal E_i'$ est $\simeq$}} \inserted{\unsure{$\mathcal E_i\simeq\mathbb Z$|$\mathcal E_i'\simeq\mathbb Z$} i.e.} $\exists!$ isom
% l. 11
\unsure{$T\xrightarrow{\sim}\mathcal E_i$|$T\xrightarrow{\sim}\mathcal E_i'$} tel que \unsure{$T\xrightarrow{k_i'}\mathcal E_i'$} s'identifie à
% l. 12
$n\operatorname{id}_T$ (de sorte qu'on trouve \unsure{$\Gamma\simeq\mathcal E_i'/n\mathcal E_i'$}
% l. 13
$\simeq T\otimes_{\mathbb Z}\mathbb Z/n\mathbb Z$, si $n=\operatorname{card}\Gamma$ ...

% l. 14
Changeant de notations, ceci revient au
% l. 15
\uline{Th}\textsuperscript{\unsure{m}} Soit $\Gamma$ un groupe fini opérant
% l. 16
fidèlt sur un groupe \uline{\unsure{extérieur|extérieurem\textsuperscript{t}}} à
% l. 17
lacets $\pi$ \inserted{de type \deleted{\illegible{1}} $(g,\nu+1)$ ($\nu\geq0$)}, en laissant fixe un $i\in I(\pi)$,
% l. 18
Alors $\Gamma$ est cyclique (si $\Gamma=\Gamma^0$) ou
% l. 19
dihédral (si $\Gamma\ne\Gamma^0$), et dans l'extension
% marge, en bas à gauche
\margin{\unsure{au} $\overline{\mathrm m}$ il y a \uline{équivalence}, si dans le th, on suppose le type $(g,\nu)$ anabélien --- sinon le th. est un peu plus général}
\mspage{190}


\scanpage{192}
\mspage{96}
% l. 01
correspondant \deleted{\illegible{2}} \unsure{$\mathcal E$} de $\Gamma$ par $\pi$,
% l. 02
Si \unsure{$\mathcal E_i$} est le normalisateur de un $L_i$ dans
% l. 03
\unsure{$\mathcal E$} [ \inserted{de sorte qu'on a} \deleted{une} \unsure{$\mathcal E_i\to\Gamma$} épi i.e. une suite exacte
% l. 04
\unsure{$1\to L_i\to\mathcal E_i\to\Gamma\to1$} ]. l'image inverse \unsure{$\mathcal E_i^0$} de
% l. 05
$\Gamma^0$ est $\simeq\mathbb Z$ donc (si $n=\operatorname{card}\Gamma^0$)\inserted{\unsure{$=[\mathcal E_i^0:L_i]$}}
% l. 06
\deleted{\unsure{$n\operatorname{id}$}} $x\to x^n$ est un iso. \deleted{et} \unsure{$\mathcal E_i^0\simeq L_i$}, qui
% l. 07
compte tenu de $k_i$ donne un
% l. 08
isom \unsure{$\mathcal E_i^0\simeq T$}, dont le composé avec
% l. 09
\deleted{\unsure{$\mathcal E_i$}} \unsure{$T\xrightarrow[\sim]{k_i}L_i\subset\mathcal E_i^0$} est $n\operatorname{id}_T$, de sorte
% l. 10
qu'on a un iso. can
% l. 11
\[\unsure{\Gamma^0\simeq\mathcal E_i^0/L_i\simeq T/nT}\]
% l. 12
(évidemment indépendant du choix de $L_i$ ...)

% l. 13
\uline{Démonstration}. Considérons l'extension
% l. 14
\unsure{$\mathcal E_i^0$} de $\Gamma^0$ par $L_i\simeq T$ ($\overset{\text{\unsure{non can}}}{\simeq}\mathbb Z$),
% l. 15
comme \unsure{$\mathcal E_i^0$|$\Gamma^0$} opère trivialem\textsuperscript{t} sur $T$ \inserted{sans torsion}, cette
% l. 16
extension (par la suite exacte de coh.
% l. 17
associée à $0\to T\to T\otimes_{\mathbb Z}\mathbb Q\to T\otimes_{\mathbb Z}\mathbb Q/\mathbb Z\to0$)
% l. 18
est canoniquem\textsuperscript{t} isomorphe à \deleted{\unsure{une}} l'extension
% l. 19
définie par un hom $\Gamma^0\to T_n=T/nT$,
% l. 20
comme image inverse de l'extension
% l. 21
$0\to T\xrightarrow{n}T\to T_n\to0$ de $T_n$ par $T$. Je dis que
\mspage{191}


\scanpage{193}
% l. 01
cet hom. est un iso (d'où \unsure{résulteraient|résulteront}
% l. 02
les autres assertions), ou ce qui revient
% l. 03
au $\overline{\mathrm m}$ puisque $\Gamma^0$ et $T_n$ sont tous
% l. 04
deux d'ordre $n$, que c'est mono.
% l. 05
Remplaçant $\Gamma^0$ par le noyau de $\Gamma^0\to\mathbb Z/n\mathbb Z$,
% l. 06
OPS que l'hom. \inserted{en question} est nul, et il faut
% l. 07
prouver que cela implique que l'action
% l. 08
de $\Gamma^0$ est triviale (ce qui, puisque par hyp.
% l. 09
l'action de $\Gamma$ est fidèle, implique $\Gamma^0=\{1\}$,
% l. 10
OK.) Donc \unsure{th} est ramené au

% l. 11
\uline{lemme fondamental}. Tout automorphisme \inserted{direct}
% l. 12
extérieur $u$ d'ordre fini \inserted{\unsure{$n$}} d'un groupe à
% l. 13
lacets $\pi$, qui fixe une classe de lacets
% l. 14
$i$ et est tel que l'extension de $\mathbb Z/n\mathbb Z$
% l. 15
par $L_i\in i$ définie par $u$ soit triviale,
% l. 16
est trivial.

% l. 17
L'hypothèse signifie que $u$ se remonte
% l. 18
en un autom $u_0$ de $\pi$, qui
% l. 19
normalise $L_i$, et qui soit aussi d'ordre \unsure{$n$},
% l. 20
(ou d'ordre fini, cela revient au \unsure{$\overline{\mathrm m}$|$\Rightarrow$} compte tenu
% l. 21
que \deleted{\unsure{les élts de}} $T$ est sans élts d'ordre fini)
\mspage{192}


\scanpage{194}
\mspage{97}
% l. 01
alors $u_0$ est trivial, i.e. cela équivaut au
% l. 02
\uline{Corollaire}. Tout autom \inserted{direct} d'ordre fini \inserted{$u_0$} d'
% l. 03
un groupe à lacets $\pi$, \inserted{qui soit direct et} qui normalise
% l. 04
un sous-groupe à lacets $L_i$ (i.e. qui
% l. 05
centralise $L_i$), est réduit à l'identité.
% marge supérieure gauche
\margin{Dans le lemme ou son \unsure{coroll.}}
\margin{\uline{NB} (le cas $g=0$, $\nu+1=1$ ou $\nu+1=2$ est trivial, le cas $\nu+1=3$ ($(g,\nu)=(0,2)$ abélien!) n'est pas trivial \unsure{par contre}, ni le cas $g=1$, $\nu+1=1$ i.e. $(g,\nu)=(1,0)$ abélien. Pourtant le résultat doit être valable aussi dans \unsure{ce|ces} cas.)}

% l. 06
Pour le démontrer, on est ramené \inserted{aussitôt} au
% l. 07
cas où $u_0$ est \deleted{d'ordre premier $p$} \inserted{tel que $u_0^p=\operatorname{id}$, où}
% l. 08
$p$ est un \unsure{nb} premier --- i.e. $u_0$ corres-
% l. 09
pond à une opération au sens strict
% l. 10
(pas seulem\textsuperscript{t} extérieure) de $\mathbb Z/p\mathbb Z$ \deleted{\unsure{par}} \unsure{à autom}
% l. 11
de $\pi$. Mais (que l'on puisse ou non
% l. 12
trouver un tel $p$) considérons le cas où
% l. 13
on sait que l'opération \uline{extérieure} de
% l. 14
$\Gamma=\mathbb Z/n\mathbb Z$ [$n=$ ordre de $u$] sur $\pi$ se réalise
% l. 15
géométriquem\textsuperscript{t}, par une opération (fidèle)
% l. 16
de $\Gamma=\mathbb Z/n\mathbb Z$ sur une $U$ de type $(g,\nu+1)$.
% l. 17
\deleted{\unsure{classes de $\pi$-conjugaison}, \illegible{1}}
% marge inférieure gauche
\margin{$U=X\setminus S$, $S\simeq I$, $X$ comp. de genre $g$ \unsure{la pointe|la partie} $s_i$ de $S$ correspondant à un pt fixe de $\Gamma$ op. sur $U$.}

% l. 18
\deleted{On sait que les (\unsure{relèvements}) des auts
% l. 19
extérieurs en \deleted{\unsure{des autom}} \inserted{une vraie opération},
% l. 20
correspondent aux pts fixes \illegible{1} de $\Gamma$
% l. 21
dans $U$; \deleted{et l'extension de} \inserted{soit $a$ le pt fixe de $\Gamma$}
% l. 22
correspondant à l'opération donnée, \unsure{de}
% l. 23
sorte qu'on a \deleted{\unsure{un}} iso. can. \unsure{compatible avec $\Gamma$}
% l. 24
$\pi=\pi_1(U,a)$.}
\mspage{193}


\scanpage{195}
% l. 01
On exprime l'hyp. de l'opération \inserted{\unsure{``véritable''}} de $\Gamma$ sur
% l. 02
$\pi$, \deleted{\illegible{1}} \unsure{centralis.} un $L_i$, en disant que
% l. 03
l'opération extérieure donne une extension
% l. 04
de $\Gamma$ par $L_i\simeq T$ triviale. Mais on sait
% l. 05
par ailleurs dans le cas géom. \inserted{et op. fidèle} qu'elle
% l. 06
n'est \uline{pas} triviale !

% l. 07
Il suffirait donc, \deleted{\illegible{1}} pour prouver le
% l. 08
lemme fondamental, de savoir que tte
% l. 09
action extérieure \inserted{[fidèle]} d'un groupe cyclique
% l. 10
sur un groupe à lacets de type \unsure{$(g,\nu+1)$}
% l. 11
est réalisable, et il suffit $\overline{\mathrm m}$ de le savoir
% l. 12
pour un groupe cyclique d'ordre
% l. 13
premier. \inserted{\unsure{est}} \unsure{Or} sauf erreur ce résultat
% l. 14
est \unsure{bien} connu ($\overline{\mathrm m}$ pour les groupes \unsure{ré-}
% l. 15
solubles ...). --- de sorte que le lemme fon-
% l. 16
damental semble démontré. J'ai seulement
% l. 17
un doute s'il est démontré dans le
% l. 18
cas général d'un $(g,\nu)$, \inserted{\deleted{\illegible{1} $g\geq2$}} ou seulement
% l. 19
pour \inserted{$g\geq2$,} $\nu=0$. Mais \deleted{à vrai dire, \illegible{2}, \inserted{cela suffit}}
% l. 20
\deleted{\illegible{2} \inserted{pour $g\geq2$}}
% marge inférieure gauche
\margin{comme th. \unsure{d'exist.} \unsure{de pts} fixes d'un groupe \illegible{1} \unsure{opér.} \unsure{de} Teichmüller}

% l. 21
\deleted{\uline{Lemme} Soit \illegible{1} une opération extérieure}
\mspage{194}


\scanpage{196}
\mspage{98}
% l. 01
\deleted{d'un groupe fini $\Gamma$ sur un \illegible{1}} s'il en
% l. 02
est bien ainsi, je pense que (pour $g\geq2$ tout
% l. 03
au moins) on s'en tire par dévissage
% l. 04
pour le cas $\nu$ quelconque --- et le cas $g=0,1$
% l. 05
demanderait \inserted{\unsure{aussi}} un traitement à part.
% l. 06
Je reviendrai là-dessus par la suite, et
% l. 07
préfère pour l'instant admettre le ``lemme
% l. 08
fondamental'', et examiner des conséquences
% l. 09
et corollaires de celui-ci.

% l. 10
Pour une action extérieure \inserted{fidèle} d'un $\Gamma$ fini
% l. 11
sur un groupe à lacets \uline{anabélien} $\pi$, \unsure{on}
% l. 12
\deleted{\unsure{a établi}} correspondant à une extension \unsure{$\mathcal E$} de $\Gamma$ par $\pi$,
% l. 13
on a donc établi a) que les scindages
% l. 14
partiels de celle-ci ne peuvent se faire
% l. 15
que sur des ss-groupes \deleted{\unsure{cycliques}} $\Gamma'$ de $\Gamma$
% l. 16
tels que $\Gamma'{}^0$ \unsure{soit} cyclique, et $\Gamma'$ dihédral
% l. 17
si $\Gamma'\ne\Gamma'{}^0$, et que b) pour tte classe
% l. 18
de \illegible{2} de tels scindages partiels
% l. 19
\unsure{sur} $\Gamma'\subset\Gamma^0$, on a un iso. can.
% l. 20
correspondant $\Gamma'\simeq T_n$ ($=T/nT$), où $n=\operatorname{card}\Gamma'$.
% marge inférieure gauche
\margin{i.e. on a établi l'existence d'une \unsure{applic.} can. \illegible{1} des \unsure{éléments} d'ordre fini de $\operatorname{Aut}_{\mathrm{lac}}(\pi)$ dans \unsure{$T\otimes\mathbb Q/T$}, satisfaisant \unsure{les} conditions \unsure{examinées} précédemment}

% l. 21
Ce sont là des résultats qu'on avait
\mspage{195}


\scanpage{197}
% l. 01
précédemment obtenus pour le cas d'une
% l. 02
opération \uline{réalisable}. [Il n'y avait pas lieu
% l. 03
d'ailleurs de se borner au cas anabélien,
% l. 04
du moment qu'on \unsure{suppose} $\pi\ne0$ \inserted{(cas essentiellement vide!)},
% l. 05
ce qui inclut les cas \inserted{abéliens} $g=0$, $\nu=2$ et
% l. 06
$g=1$, $\nu=0$, pour lesquels un traitement
% l. 07
direct est possible, et a déjà été donné
% l. 08
essentiellement, les cas-là \deleted{\unsure{étant} réalisables,} \inserted{\unsure{où|si l'on}}
% l. 09
part d'une \uline{extension} (pas d'une action
% l. 10
``extérieure'' i.e. ici d'une action tout court
% l. 11
de $\Gamma$ sur $\pi\simeq\mathbb Z$ ou $\mathbb Z^2$), \unsure{étant} tjrs \unsure{réalisables}.
% l. 12
(\uline{NB} \deleted{\unsure{l'hyp.}} \inserted{Dans ce} cas, l'hypothèse d'une
% l. 13
action fidèle est remplacée par celle que
% l. 14
l'extension n'est \deleted{\unsure{pas scindable}} \inserted{une extension \uline{produit}} \unsure{sur|sous} aucun
% l. 15
sous-groupe $\Gamma'\subset\Gamma^0$ avec $\Gamma'\ne1$ ....).

% l. 16
Il reste cependant d'autres résultats du
% l. 17
cas réalisable, qu'il faudrait examiner
% l. 18
dans le cas général :

% l. 19
c) Si $\Gamma$ est un sous-groupe \deleted{\unsure{discret}} fini
% l. 20
de $\operatorname{Aut}_{\mathrm{lac}}(\pi)$ i.e. un grpe fini opérant \uline{\unsure{fidèlt}}
% l. 21
sur $\pi$, a-t-on $\pi^\Gamma=\{1\}$ ?
% marge inférieure gauche
\margin{OPS $\Gamma=\Gamma^0$ donc $\Gamma$ engendré par un autom. direct \unsure{$u$} ou $\Gamma$ engendré par un \unsure{anti-autom.} d'ordre 2}
\mspage{196}


\scanpage{198}
\mspage{99}
% l. 01
d) \unsure{Si} $\Gamma',\Gamma''$ sont des ss-groupes finis de \unsure{$\mathcal E$}
% l. 02
\unsure{(ext.} du groupe fini $\Gamma$ par $\pi$, correspondant
% l. 03
à une opération \deleted{\unsure{respect}} extérieure \inserted{(respectant l'orientation)} fidèle), \deleted{\illegible{1}}
% l. 04
tels que $\Gamma'\cap\Gamma''\ne(1)$, alors \deleted{l'un des
% l. 05
deux contient-il l'autre} ?
\inserted{$\Gamma'$ et $\Gamma''$ sont-ils contenus dans le $\pi$-conjugué d'un ss-gpe fini $\Gamma'''$ de \unsure{$\mathcal E$} [i.e. tt sous-\unsure{groupe} section $\ne1$ est contenu dans un unique sous-groupe section \unsure{maximal}, mod. $\pi$-conjugaison}
% marge supérieure gauche
\margin{\uuline{NB} Deviendrait faux en \deleted{\illegible{1}} \unsure{plaçant} dans le groupe $\operatorname{Aut}_{\mathrm{lac}}(\pi)$ \illegible{2} \unsure{l'un en ss-gr.} \illegible{2} \unsure{l'un des deux ss.} \unsure{contenu dans} \illegible{4}}

% l. 06
e) Pour tt \unsure{ss-gpe}, \illegible{2} \deleted{\illegible{1}} $\Gamma'$ de $\Gamma$,
% l. 07
\unsure{l'ens.} des classes \deleted{\unsure{d'extensions}} $\pi$-conjugaison
% l. 08
de relèvements de $\Gamma'$ en \unsure{$\mathcal E$} est-il fini ?
% l. 09
[OPS $\Gamma'=\Gamma$, et $\Gamma$ cyclique (si $\Gamma=\Gamma^0$) ou dihédral
% l. 10
(sinon)]
% l. 11
\inserted{Dans le cas \unsure{c)|e)}, OPS $\Gamma$ cyclique d'ordre premier, et
% l. 12
c'est OK s'il est acquis qu'une \unsure{op. extérieure} \illegible{3}
% l. 13
\unsure{que} \illegible{4} \unsure{$\pi_{g,\nu}$} \deleted{\illegible{3}} \unsure{d).}} \deleted{En appliquant}
% marge gauche, au milieu
\margin{\unsure{De} $\overline{\mathrm m}$ e) est établi si \unsure{on} admet que les op. ext. de \unsure{groupes cycliques} \unsure{sur un} $\pi\simeq\pi_{g,\nu}$ sont réalisables}

% l. 14
\deleted{Ex. examinons \illegible{2}
% l. 15
la \unsure{joncture|jonction} ``forage de trous'', $\Gamma'$ et $\Gamma''$
% l. 16
opérant extérieurement sur un $(\pi';i)$
% l. 17
$i\in I(\pi')$ --- leur action sur $\pi$ est
% l. 18
déduite de celle sur $(\pi',i)$ en ``bouchant
% l. 19
le trou $i$''. Ceci dit, $\Gamma'$ et $\Gamma''$ sont identifiés
% l. 20
aux ss-groupes finis de $\operatorname{Autext}_{\mathrm{lac}}(\pi';i)$
% l. 21
On a mieux :
% l. 22
\unsure{$\Gamma=\Gamma'\cap\Gamma''\ne(1)$}}

% Un trait oblique part du coin inférieur gauche du cadre et traverse « Soit » et « et considérons » ; sa portée comme biffure est incertaine.
\unsure{\deleted{
% l. 23
\uline{Théorème}. Soit $\Gamma_0$ \deleted{\illegible{1}} un sous-groupe fini de
% l. 24
$A=\operatorname{Aut}_{\mathrm{lac}}(\pi)^0$, \inserted{$\Gamma_0\ne1$}, \deleted{\illegible{1} et considérons} Alors le normalis. $N$
% l. 25
de \unsure{$\Gamma$|$\Gamma_0$} dans $A$ est égal à son centralisateur \unsure{$Z$|$\mathcal Z$}
}}
\margin{\uuline{Ultra fort!} \unsure{Si $g=0$}}
\mspage{197}


\scanpage{199}
% Un trait oblique traverse le corps de cette page ; sa portée comme biffure est incertaine.
\unsure{\deleted{
% l. 01
et l'ensemble des élts d'ordre fini
% l. 02
de \unsure{$Z$|$\mathcal Z$} est un groupe cyclique \inserted{\unsure{$F$|$\Gamma$}}, \deleted{contenu
% l. 03
dans} qui est aussi le plus grand sous-
% l. 04
groupe fini de $A$ contenant $\Gamma_0$, et est
% l. 05
contenu dans le centre de \unsure{$Z$|$\mathcal Z$}.

% l. 06
Le fait que le normalisateur $N$
% l. 07
centralise $\Gamma_0$ provient du fait que l'on
% l. 08
a un iso can. (compatible aux \unsure{transports}
% l. 09
de \unsure{structures})
% l. 10
\[\Gamma_0\xrightarrow{\sim}T_{n_0}\qquad(\text{si }n_0=\operatorname{card}\Gamma_0)\]
% l. 11
(NB si $g\in\operatorname{Aut}_{\mathrm{lac}}(\pi)$ renversant l'orientation,
% l. 12
et \unsure{normalise|centralise} $\Gamma_0$, alors \deleted{\unsure{$g$}} $(\operatorname{int}g)$ induit
% l. 13
la symétrie $\sigma\to\sigma^{-1}$ sur $\Gamma_0$, par le $\overline{\mathrm m}$
% l. 14
argument) \deleted{\illegible{1}} Notons que si
% l. 15
$u$ est un élt d'ordre fini de \unsure{$Z$|$\mathcal Z$}, de
% l. 16
sorte que $\Gamma_0,u$ engendrent un ss-grpe
% l. 17
comm. fini de \unsure{$Z$|$\mathcal Z$}, \deleted{on voit que} \inserted{comme on sait}
% l. 18
qu'un ss-groupe fini de $A$ est cyclique,
% l. 19
il s'ensuit que $u$ est contenu dans
% l. 20
un sous-groupe cyclique de \unsure{$Z$|$\mathcal Z$} \unsure{cont\textsuperscript{t}} $\Gamma_0$,
% l. 21
\deleted{\unsure{donc}} soit $\Gamma'$. choisissons maintenant
% l. 22
une réalisation géom. de l'action \inserted{ext.} de
% l. 23
$\Gamma_0$ sur \deleted{\illegible{1}} $\pi$, soit $(X,S)$ avec $U=X-S$ ;
}}
\mspage{198}


\scanpage{200}
\mspage{100}
% Un long trait oblique traverse le texte ; sa portée comme biffure est incertaine.
\unsure{\deleted{
% l. 01
on sait que le relèvement de l'action
% l. 02
extérieure de $\Gamma_0$ en une action est
% l. 03
décrit par un unique pt $a\in U$.
% l. 04
\deleted{Soit} Je dis que $\Gamma'$ \deleted{est réalisable} (ou son
% l. 05
action \deleted{extérieure} sur $\pi\simeq\pi_1(U,a)$) est réalisable
% l. 06
de \uuline{façon unique} par une action de $\Gamma'$
\margin{\uuline{pas vrai!}}
\inserted{et induisant l'action déjà donnée de \unsure{$\Gamma$|$\Gamma_0$}}
% l. 07
sur $(X,S)$), fixant $a$, le fait qu'il
% l. 08
soit \inserted{bien} réalisable ainsi provient du fait
% l. 09
(que nous admettons) que tte action
% l. 10
extérieure d'un groupe cyclique \inserted{(ici $\Gamma'$)} sur un
% l. 11
groupe à lacets \inserted{$\pi$} de type $(g,\nu)$ est réalisable,
% l. 12
et de l'unicité d'une telle réalisation
% l. 13
(dans le cas de $\Gamma_0$). Pour l'\uline{unicité},
% l. 14
cela \deleted{\illegible{1}} revient au

% l. 15
\uline{Lemme}. Soit $u_0$ un autom d'ordre fini $n_0\ne1$
% l. 16
d'une surface \inserted{orientable} \deleted{\unsure{supposée} $X^{u_0}\ne\emptyset$.} \inserted{$U=X-S$, $X$ (compacte, \deleted{\illegible{1}}, et} laissant fixe \inserted{un pt $a\in U$} \illegible{1}
% l. 17
soit $d\in\mathbb N^*$, (S'il existe un automorphisme \unsure{$u$} de
% l. 18
\deleted{\illegible{1}} \inserted{$(X,S)$} tel que $u^d=u_0$, \inserted{$u(a)=a$}, alors $u$ est uniquement
\margin{\uline{faux!}}
% l. 19
déterminé par son \deleted{\unsure{classe d'}} action sur $\pi_1(U,a)$;
% l. 20
et en tout état de cause le groupe cyclique
% l. 21
d'autom. de $(X,S)$ engendré par $u$ est \unsure{uniquem\textsuperscript{t}}
% l. 22
\inserted{déterminé (il est d'ordre $nd$).}
}}
\mspage{199}
